\begin{abstract}
We prove Donovan's conjecture over every algebraically closed field
\(K\) of characteristic \(p>0\). For each fixed \(K\) and bound on
defect-group order, blocks of finite groups represent only finitely
many \(K\)-linear Morita equivalence classes. The defect groups need
not be abelian, and the result includes \(p=2\).
\end{abstract}

\section{Introduction}
\label{sec:introduction}

Let \(p\) be a prime and let \(K\) be an algebraically closed field
of characteristic \(p\).
A block of a finite group algebra \(KG\) is a summand \(KGb\)
defined by a primitive central idempotent \(b\).  Its defect group
is a \(p\)-subgroup of \(G\), determined up to conjugacy, that
measures how far the block is from being semisimple.  Two blocks
are \(K\)-linearly Morita equivalent when their module categories
are equivalent by a \(K\)-linear functor.  Donovan's conjecture
asks whether a fixed finite \(p\)-group can occur as the defect
group of only finitely many blocks up to this equivalence.
The question constrains an entire representation category by
local group data: neither the order of \(G\) nor the dimensions
of its simple modules are fixed, and the block need not be principal.
We prove the conjecture in the following bounded-order form.

\begin{theorem}[Donovan's conjecture]
\label{thm:donovan}
Fix a prime \(p\), an algebraically closed field \(K\) of
characteristic \(p\), and a positive integer \(M\).
As \(G\) ranges over finite groups and \(b\) over block idempotents of
\(KG\) with defect groups of order at most \(M\), only finitely many
\(K\)-linear Morita equivalence classes of \(KGb\) occur.
In particular, over this field \(K\), Donovan's conjecture holds for every finite
\(p\)-group, including nonabelian groups.
\end{theorem}

The bounded-order formulation is equivalent to the formulation with
one fixed defect group: there are only finitely many isomorphism
types of groups of bounded order.  The field \(K\) is fixed while
the finite group and block vary; the theorem concerns equivalences
of module categories over that field.  The resulting finite list also
bounds Cartan entries, Loewy lengths of basic algebras, and, in each
fixed degree, dimensions of extension groups between simple modules.

The proof first treats \(k=\overline{\mathbb F}_p\), which is the
coefficient field used in the main argument below.  After choosing an
embedding \(k\hookrightarrow K\), every block over \(K\) descends to a
unique block over \(k\) with the same defect groups, by the coefficient
comparison of \cite[Section~2.1]{BlockwiseAlperinWeight}.  Extending the
Morita bimodules then transfers the finite list from \(k\) to \(K\);
this final step is given in Section~\ref{sec:coefficient-extension}.

\subsection{Previous work}

The conjecture is attributed to Peter Donovan and appears as
Conjecture~M in Alperin's account of problems in group representation
theory \cite{Alperin1980}.
Its Cartan-boundedness component grows out of a question of Brauer;
Hiss emphasized the additional rationality obstruction
\cite{HissBrauer}.  Kessar proved that field-valued Donovan
finiteness separates into Cartan boundedness and bounded
Morita--Frobenius numbers \cite{KessarRemark}.

Scopes proved finiteness for symmetric-group blocks, and Kessar
treated their double-cover families \cite{Scopes,KessarSymmetric}.
For unipotent blocks of general linear groups, Jost's Morita
reductions bound the rank in terms of the defect \cite{Jost}.
Combined with rationality, this gives
finiteness across both ranks and field sizes; see
\cite[Example~5.3(2)--(3)]{HissBrauer} for these results and their
relationship.  Thus uniformity in these two parameters already has
an important type-A precedent.  Hiss--Kessar developed Scopes
reductions for unipotent blocks of classical groups, including
orthogonal blocks with degenerate-symbol characters in the sequel
\cite{HissKessarScopes,HissKessarScopesII}.

Eaton--Livesey proved the field-valued conjecture for abelian
\(2\)-groups \cite{EatonLivesey}.  Eaton--Eisele--Livesey established
an integral finiteness criterion and an abelian-defect reduction to
quasisimple blocks; together with Farrell--Kessar's rationality
bounds, this leaves Cartan boundedness as the remaining condition
in that reduction \cite{EEL,FarrellKessar}.
An--Eaton treated extraspecial groups of order \(p^3\) and exponent
\(p\) for \(p\ge5\), and reduced the corresponding case at \(p=3\)
to a quasisimple Cartan bound \cite{AnEaton}.
Eaton treated blocks with Suzuki \(2\)-group defects
\cite{EatonSuzuki}.  The August 2026 preprint
\cite[Theorem~1.1 and Corollary~1.2]{Quaternion2026} proves
integral finiteness for quaternion defect groups and deduces
finiteness for all tame blocks over \(k\).

For extensions, K\"ulshammer introduced crossed products and
algebraic-group actions into the reduction of Donovan's conjecture
\cite{Kulshammer1995}.  Eisele developed integral crossed-product
reductions and proved finiteness of integral block Picard groups
\cite{EiseleReduction,EiselePicard}.  Their prime-to-\(p\)
crossed-product finiteness is an antecedent of the extension method
below; the additional comparison here retains Sylow \(p\)-actions
and their specified inner multiplication factors.

\subsection{The two uniformity problems}

The general theorem requires both multiplicity control and control
of field of definition.  Bounds on Cartan entries alone do not
supply the latter.  A Cartan entry records a composition multiplicity
in a projective cover.  Their sum is the dimension of a basic algebra,
which represents the Morita class with all simple modules
one-dimensional.  Bounded dimensions do not by themselves confine
its structure constants to one finite field.  The Brauer--Feit
bound on ordinary characters also bounds the number of simple
modules in a block of bounded defect \cite{BrauerFeit}; thus a
Cartan-entry bound gives the required dimension bound.  If both
the dimension
and the size of a finite field of definition are bounded, only finitely many
multiplication tables can occur.

The multiplicity problem must also be uniform in parameters that
the defect does not bound.  In a group of Lie type, finite tori can
grow even when the chosen block has small defect; a bound for a
principal block with small Sylow subgroup does not address this
case.  The geometric construction below counts indecomposable
projective factors rather than their vector-space dimensions,
and estimates characters after projection to the chosen block.
Classical rank requires a separate reduction to bounded rank or
to a cuspidal series whose Harish--Chandra weight is bounded in
terms of the defect order.

Likewise, passing from quasisimple groups to
arbitrary groups needs more than a non-equivariant Morita--Frobenius
bound on each component.  Here coefficient Frobenius raises scalars
to their \(p\)th powers, and a Morita--Frobenius bound controls how
many iterations are needed to return to the Morita class.
A group extension also specifies how its homogeneous components
multiply: choices of coset representatives have products differing
by elements of the normal subgroup.  These inner factors remain
part of the crossed algebra throughout the proof.  We construct
Frobenius comparisons compatible with the actions of Sylow
\(p\)-subgroups in the crossed extensions and with those factors.
Scalar discrepancies can
then be removed; arbitrary central-unit discrepancies would not
permit this normalization.

\subsection{The proof and its new ingredients}

The proof first establishes a uniform Cartan bound for quasisimple
blocks of bounded defect.  It then reduces a general block to crossed
extensions of finitely many integral base orders and proves
finiteness of the remaining crossed data.  The Cartan estimate bounds the sizes of the basic algebras;
the comparison controls their multiplication and descent data.
Four constructions address the parameters that can still be unbounded.  The two main outputs
are Theorem~\ref{thm:quasisimple-cartan}, which supplies the Cartan
bound, and Proposition~\ref{prop:sylow-comparison}, which supplies
the equivariant comparison.  Figure~\ref{fig:proof-interfaces}
shows their distinct roles in the final extension argument.

\paragraph{Projective multiplicities in fixed root data.}
Section~\ref{sec:geometry} constructs projective characters from
tilting objects on the two-flag space, following Eteve's construction
\cite{Eteve}.  The essential estimate bounds
the sum of projective multiplicities in extraordinary stalks.
Regular Lang local systems and a Kummer-cohomology calculation remove
dependence on the size of the finite torus \(p\)-part.
An ordinary-coordinate projection then bounds the character norms
of projective indecomposable modules (PIMs).
This estimate is uniform in the field size and monodromy parameter,
but its root data remain fixed.

\paragraph{Classical rank and projective cones.}
Sections~\ref{sec:families}--\ref{sec:runners} deal with classical
groups of unbounded rank.  A single-cycle Jordan-label calculation
makes the required ordinary induction rules available beyond uniform
class functions.  Harish--Chandra moves then act on projected cones
of projective characters.  Most moves are exact permutations after
normalization.  Here the cones are the nonnegative spans of PIM
characters in selected ordinary-character coordinates.
The remaining moves have summable errors, including
a two-edge detour for the exceptional cohook configuration.
Determinant and cone-volume estimates convert this control into
bounds on the PIM columns.  These moves need not be Morita
equivalences.  The runner language belongs to the lineage of
Scopes's abacus methods, Jost's general-linear reductions, and
Hiss--Kessar's classical reductions
\cite{Scopes,Jost,HissKessarScopes,HissKessarScopesII}.
Here the transported object is a cone
of projective characters; its positivity and determinant bounds
carry a numerical estimate through all classical families.

\paragraph{The cuspidal endpoint.}
Runner reductions can leave a cuspidal label of unbounded rank:
a staircase partition in the unitary case, or a two-row symbol
with each row an initial interval of nonnegative integers in the
other classical types.
We call these labels cuspidal triangles.
Section~\ref{sec:endpoint} treats this endpoint directly by modular
Harish--Chandra induction.  We construct the intertwiners, remove
their extension obstruction, count the resulting Hecke-algebra
simple types, and prove that their heads cover the required simple
modules.  The bound depends on the Harish--Chandra weight above the
triangle, not on the triangle's rank.

\paragraph{Equivariant descent of crossed extensions.}
Sections~\ref{sec:extensions} and~\ref{sec:periodicity} address
arbitrary finite groups.  The generalized Fitting reduction leads
to base blocks equipped with an actual inner factor system.
We construct integral Morita comparisons with bounded Frobenius
period whose reductions respect the Sylow actions and the
specified inner factors.
Their multiplication error is scalar, which can be removed over
\(k\); an arbitrary central-unit error would not suffice.
Finite integral Picard groups then turn these comparisons into
Frobenius descent retaining the chosen base algebra.
Lang's theorem and a Sylow restriction
argument leave only finitely many crossed systems.

The role of the integral theory is confined to the base orders and
their comparisons.  The final removal of the large \(p'\)-kernel
and the final crossed-product finiteness are over \(k\).
The finite lists of integral identity-component orders and the
integral Morita bimodules with compatible projective transports
also serve the subsequent integral companion over \(W(k)\)
\cite{OpenAIIntegralDonovan2026}.  The field proof given here is
independent of that companion.
The proof invokes established structural, block-theoretic and
finite-reductive-group theorems with their hypotheses specified
where they are used.  The new uniformity, runner, endpoint and
equivariance arguments are proved below.

\begin{figure}[tbp]
\centering
\setlength{\fboxsep}{6pt}
\begin{tabular}{@{}c@{\hspace{1em}}c@{}}
\fbox{\begin{minipage}{.40\linewidth}\centering\small
Fixed-root geometry; classical labels,\par
runner cones and triangle endpoints\par
Sections~\ref{sec:geometry}--\ref{sec:endpoint}
\end{minipage}}
&
\fbox{\begin{minipage}{.40\linewidth}\centering\small
Frobenius comparisons compatible\par
with Sylow actions and inner factors\par
Section~\ref{sec:periodicity}
\end{minipage}}
\\[3pt]
$\Downarrow$ & $\Downarrow$ \\[3pt]
\fbox{\begin{minipage}{.40\linewidth}\centering\small
Bounded quasisimple Cartan entries\par
Theorem~\ref{thm:quasisimple-cartan}
\end{minipage}}
&
\fbox{\begin{minipage}{.40\linewidth}\centering\small
Bounded equivariant Frobenius period\par
Proposition~\ref{prop:sylow-comparison}
\end{minipage}}
\\[3pt]
\multicolumn{2}{c}{$\searrow\qquad\qquad\swarrow$}\\[3pt]
\multicolumn{2}{c}{\fbox{\begin{minipage}{.83\linewidth}\centering\small
Generalized Fitting reduction and integral base finiteness\par
Comparison with trivial action: finitely many integral base orders\par
Full Sylow comparison: finitely many Sylow crossed systems\par
with the base algebra fixed\par
Section~\ref{sec:extensions}
\end{minipage}}}\\[3pt]
\multicolumn{2}{c}{$\Downarrow$}\\[3pt]
\multicolumn{2}{c}{\fbox{\begin{minipage}{.83\linewidth}\centering\small
Removal of the large $p'$-kernel and Sylow restriction\par
Finitely many $k$-linear Morita classes for bounded defect order\par
Theorem~\ref{thm:extension-finiteness}
\end{minipage}}}
\end{tabular}
\caption{The two controls used in the proof.  Arrows indicate uses
of the displayed outputs together with the structural and finiteness
inputs stated in Section~\ref{sec:extensions}.  The trivial-subgroup
case of the comparison supplies integral base finiteness; its full
Sylow form controls the multiplication in crossed extensions while
keeping the specified base algebra fixed.  Here
\(k=\overline{\mathbb F}_p\); Section~\ref{sec:coefficient-extension}
extends the resulting finite list to each fixed algebraically closed
field \(K\) of characteristic \(p\), completing
Theorem~\ref{thm:donovan}.}
\label{fig:proof-interfaces}
\end{figure}


Section~\ref{sec:preliminaries} fixes the coefficient conventions
and proves the transfer estimates used throughout.  The first
milestone is assembled in Section~\ref{sec:families}, using the
fixed-root estimate of Section~\ref{sec:geometry} and the classical
arguments developed in Sections~\ref{sec:labels}--\ref{sec:endpoint}.
Section~\ref{sec:extensions} states the precise comparison property
and proves that it suffices for finiteness.
Section~\ref{sec:periodicity} constructs that comparison with the
specified multiplication factors and proves finiteness over
\(k=\overline{\mathbb F}_p\).  Section~\ref{sec:coefficient-extension}
then transfers the finite list to the fixed algebraically closed field
\(K\), completing the proof of Theorem~\ref{thm:donovan}.
